\documentclass[conference]{IEEEtran}

\usepackage{cite}
\usepackage{amsmath,amssymb,amsfonts}
\usepackage{algorithmic}
\usepackage{graphicx}
\usepackage{textcomp}
\usepackage{xcolor}
\usepackage{url}

\begin{document}

\title{Journal Finder for Computer Science Articles}

\author{
\IEEEauthorblockN{Data Mining Final Project}
\IEEEauthorblockA{
Department of Computer Engineering\\
}
}

\maketitle

\begin{abstract}
Selecting a suitable journal is a common challenge for researchers because journal scope, subject area, and article terminology must align. This project builds a journal finder for computer science articles using a publication database containing abstracts, journals, keywords, and subject categories. The system recommends the five most relevant journals for a new article abstract and generates topic clusters for the corpus.
\end{abstract}

\begin{IEEEkeywords}
journal recommendation, text mining, TF-IDF, cosine similarity, clustering, computer science publications
\end{IEEEkeywords}

\section{Introduction}
Journal articles are one of the most important research outputs. For authors, choosing a journal that fits the content of an article can improve the chance of reaching the right audience. The goal of this project is to use data mining methods to recommend relevant computer science journals based on an article abstract. The project also generates topic clusters from the same corpus so that subject areas in the publication database can be inspected at a high level.

\section{Related Work}
Text mining is widely used in document retrieval, topic discovery, and recommendation systems. Vector space models such as TF-IDF represent documents by term importance and support efficient similarity matching. Salton and Buckley showed that term weighting is a practical foundation for information retrieval because it balances term frequency with collection-level rarity \cite{salton1988term}. In this project, the same idea is used to represent each article as a weighted vector.

Cosine similarity is a standard method for ranking documents by lexical and topical overlap in sparse vector spaces. It is especially suitable for TF-IDF because the angle between vectors compares term distribution while reducing the effect of document length. For journal recommendation, this makes it possible to compare a new abstract against thousands of existing article records.

Clustering methods are also commonly used to discover groups of similar documents and summarize large text collections by topic. KMeans is a simple and interpretable baseline for partitioning vectorized text into topical groups \cite{macqueen1967methods}. Sebastiani's survey of text categorization also supports the use of machine learning methods for organizing and classifying textual documents \cite{sebastiani2002machine}. These studies motivate the use of TF-IDF, cosine similarity, and KMeans as transparent data mining methods for the assignment.

\section{Dataset}
The provided SQLite database contains publication records from computer science journals. The main entities used in this project are AcademicRecord for article metadata, AcademicRecordAbstract for abstract text, Publication for journal names, AcademicRecordKeyword for author keywords, AcademicRecordKeywordPlus for Web of Science keyword terms, and AcademicRecordSubject for assigned subject categories.

After joining these entities and removing records without usable abstracts, the final dataset contains 23,061 articles, 455 journals, and 80 distinct subject areas. Each article record is represented by its title, abstract, author keywords, keyword plus terms, and subject labels.

\section{Methodology}
The preprocessing step removes HTML tags from abstracts, normalizes whitespace, lowercases text, and combines title, abstract, keywords, keyword plus terms, and subjects into one training document for each article. This combined text gives the model access to both the article content and the subject metadata provided by the database.

For journal recommendation, the system fits a TF-IDF vectorizer on all article training documents. The vectorizer uses English stop-word filtering, unigram and bigram features, document frequency thresholds, sublinear term frequency, L2 normalization, and a maximum of 40,000 features. When a user enters a new abstract, the abstract is cleaned and transformed with the same vectorizer. Cosine similarity is calculated between the input abstract and known articles. The most similar articles are grouped by journal, and journal scores are calculated from the best and average similarity values. The final output is the top five journals.

For topic discovery, the same TF-IDF representation is clustered using KMeans. Each cluster is summarized by top centroid terms, dominant subject categories, and sample journals. This provides a compact overview of the major computer science topics represented in the database.

\section{Implementation}
The implementation is organized as reusable Python modules under the \texttt{src} directory, a Jupyter Notebook for reproducible analysis, and a Streamlit application for interactive use. The data loader module reads and joins the SQLite tables. The preprocessing module normalizes text fields. The modeling module trains TF-IDF representations and produces topic clusters. The recommender module ranks journals from an input abstract.

The Streamlit interface accepts an abstract, displays recommended journals, and shows generated topic clusters. The default recommendation behavior returns the top five journals, which matches the project specification.

\section{Results and Evaluation}
The final dataset used by the system contains 23,061 articles with abstracts, 455 journals, and 80 distinct subject areas. The TF-IDF representation contains up to 40,000 features after applying document frequency thresholds and English stop-word filtering.

For a sample abstract about machine learning methods for software defect prediction, the full-dataset system returned these five relevant journals:

\begin{enumerate}
\item JOURNAL OF WEB ENGINEERING
\item AUTOMATED SOFTWARE ENGINEERING
\item IEEE TRANSACTIONS ON SOFTWARE ENGINEERING
\item EMPIRICAL SOFTWARE ENGINEERING
\item SOFTWARE QUALITY JOURNAL
\end{enumerate}

The topic clustering stage generated interpretable clusters with representative terms such as software, artificial intelligence, theory and methods, architecture, hardware, cloud computing, telecommunications, networks, wireless systems, optimization, computational biology, and library science.

The project is evaluated with functional checks. The tests verify that database loading returns article records with journals and abstracts, text preprocessing removes HTML tags and produces searchable plain text, recommendation returns five journals for a sufficiently long abstract, scores are sorted in descending order, and clustering produces topic groups with representative terms and subjects. The Jupyter Notebook executes from start to finish and stores output cells, and the Streamlit interface is available as the project software component.

\section{Conclusion}
The project demonstrates a practical journal finder tailored to computer science subject areas. TF-IDF and cosine similarity provide a transparent baseline for journal recommendation, while KMeans clustering gives an overview of major topics in the dataset. Future work can compare this baseline with transformer embeddings, supervised journal classification, and citation-aware ranking.

\begin{thebibliography}{00}
\bibitem{salton1988term}
G. Salton and C. Buckley, ``Term-weighting approaches in automatic text retrieval,'' \emph{Information Processing \& Management}, vol. 24, no. 5, pp. 513--523, 1988.

\bibitem{macqueen1967methods}
J. MacQueen, ``Some methods for classification and analysis of multivariate observations,'' in \emph{Proceedings of the Fifth Berkeley Symposium on Mathematical Statistics and Probability}, 1967, pp. 281--297.

\bibitem{sebastiani2002machine}
F. Sebastiani, ``Machine learning in automated text categorization,'' \emph{ACM Computing Surveys}, vol. 34, no. 1, pp. 1--47, 2002.
\end{thebibliography}

\end{document}
